\section{Model}

\begin{frame}
	\frametitle{Causal TCN StepNet Architecture}
	\begin{columns}[t]
		\begin{column}{0.48\textwidth}
			\begin{block}{Backbone}
				\begin{itemize}
					\item Input projection: 4 channels $\rightarrow$ 64 channels
					\item 4 causal residual convolution blocks
					\item Kernel size 5, dilations $1,2,4,8$
					\item BatchNorm, GELU, dropout 0.15
				\end{itemize}
			\end{block}
		\end{column}
		\begin{column}{0.48\textwidth}
			\begin{block}{Two heads}
				\begin{itemize}
					\item Event head: per-sample step probability
					\item Count head: expected number of steps in the window
					\item Causal convolutions support streaming inference
				\end{itemize}
			\end{block}
		\end{column}
	\end{columns}

	\vspace{0.2cm}
	\begin{center}
		\small
		Acc window $\rightarrow$ feature normalization $\rightarrow$ causal TCN $\rightarrow$ event logits + count estimate
	\end{center}
\end{frame}

\begin{frame}
	\frametitle{Models Compared in the Report}
	\scriptsize
	\begin{table}
		\centering
		\begin{tabular}{p{0.22\textwidth}p{0.28\textwidth}p{0.33\textwidth}}
			\toprule
			Method & Strengths & Observed limitations \\
			\midrule
			ACF baseline & No training, lightweight, interpretable cadence estimate & Needs quasi-periodic walking; weaker at starts/stops and irregular arm motion \\
			LSTM StepNet & Learns temporal patterns and gait rhythm & Sequential recurrence is slower to train/infer; no final checkpoint retained \\
			Causal TCN StepNet & Parallel training, fixed streaming context, event + count heads & Needs labeled data and threshold calibration \\
			\bottomrule
		\end{tabular}
	\end{table}

	\begin{block}{Final choice}
		TCN was selected for the accuracy target; ACF remains a strong low-cost fallback. Final TCN test count MAE: 0.484 steps/window.
	\end{block}
\end{frame}

\begin{frame}
	\frametitle{ACF Baseline: Classical Signal Pipeline}
	\footnotesize
	\begin{columns}[t]
		\begin{column}{0.50\textwidth}
			\begin{block}{Processing steps}
				\begin{enumerate}
					\item Compute $a_{mag}(t)$ from tri-axial acceleration
					\item Remove slow baseline drift and smooth residual motion
					\item Select strongest autocorrelation peak in 1.2--2.8 Hz walking band
					\item Accept local maxima above adaptive threshold and separated by $\rho P$ samples
				\end{enumerate}
			\end{block}
		\end{column}
		\begin{column}{0.46\textwidth}
			\begin{exampleblock}{When it works}
				Regular walking creates a stable autocorrelation peak, so period and peak-distance estimates are reliable.
			\end{exampleblock}
			\begin{alertblock}{Failure modes}
				Start/stop transitions, irregular arm swing, and non-walking periodic motion weaken stationarity.
			\end{alertblock}
		\end{column}
	\end{columns}

	\begin{center}
		\scriptsize Self-recorded ACF results: test1 $84\rightarrow88$, test2 $100\rightarrow97$; MAE 3.5 steps, MAPE 3.88\%.
	\end{center}
\end{frame}